\documentclass[11pt,letterpaper]{article}
\usepackage[T1]{fontenc}
\usepackage{lmodern,microtype}
\usepackage[margin=.7in,headheight=15pt]{geometry}
\usepackage{amsmath,amssymb,mathtools}
\usepackage{xcolor,booktabs,array,enumitem,fancyhdr,hyperref}

\definecolor{ink}{HTML}{222834}
\definecolor{indigo}{HTML}{25355D}
\definecolor{muted}{HTML}{505A6B}
\definecolor{keyink}{HTML}{205D52}
\color{ink}
\hypersetup{colorlinks=true,linkcolor=indigo,urlcolor=indigo,
  pdftitle={Expectation, Variance, and Discrete Distributions: 50-minute Quiz},
  pdfauthor={Chance, Data, and Decisions}}
\setlength{\parindent}{0pt}
\setlength{\parskip}{4pt}
\setlength{\abovedisplayskip}{7pt}
\setlength{\belowdisplayskip}{7pt}
\clubpenalty=10000
\widowpenalty=10000
\setlist[enumerate]{label=(\alph*),leftmargin=1.7em,
  itemsep=4pt,topsep=5pt,parsep=2pt}

\newcommand{\Prob}{\mathbb P}
\newcommand{\E}{\mathbb E}
\DeclareMathOperator{\Var}{Var}
\newcommand{\edition}{\ifanswers Instructor key\else Student quiz\fi}
\newcommand{\partpts}[1]{\unskip\nobreak\hspace{.3em}%
  \mbox{\small\color{muted}(#1 points)}}
\newcommand{\credit}[1]{\par{\footnotesize\color{muted}\textit{Credit:} #1}}
\newcommand{\problem}[4]{\par\addvspace{8pt}%
  {\large\bfseries\color{indigo}#1.\ #2}\hfill
  {\small\color{muted}\textbf{#3 points} $\cdot$ about #4 minutes}\par\nobreak}
\NewDocumentEnvironment{solution}{O{.8in}+b}{%
  \ifanswers
    \par\begingroup\small\color{keyink}\textbf{Solution.} #2\par\endgroup
  \else
    \par\nobreak\vspace*{#1}
  \fi
}{}

\pagestyle{fancy}
\fancyhf{}
\fancyhead[L]{\small\color{muted}Expectation and distributions quiz}
\fancyhead[R]{\small\color{muted}Chance, Data, and Decisions}
\fancyfoot[L]{\footnotesize\color{muted}\edition}
\fancyfoot[R]{\footnotesize\color{muted}Page \thepage\ of \pageref*{quiz-last-page}}
\renewcommand{\headrulewidth}{.3pt}

\begin{document}

{\LARGE\bfseries\color{indigo}Expectation, Variance,\\[2pt]
and Discrete Distributions}\par
{\small\color{muted}50 minutes \quad $\bullet$ \quad 50 points
  \quad $\bullet$ \quad \edition}\par
\ifanswers
  \textbf{Grading note.} Award method credit for correct reasoning after an
  earlier error; avoid repeated deductions for the same arithmetic mistake.
  Suggested working time is 46 minutes, leaving 4 minutes for checking.\par
\else
  \smallskip
  \textbf{Name:}\ \rule{.49\linewidth}{.4pt}\hfill
  \textbf{Date:}\ \rule{.20\linewidth}{.4pt}\par
\fi
\textbf{Instructions.} Answer all four problems and show your reasoning.
No calculator is needed; exact fractions or exact expressions with powers
and binomial coefficients are acceptable. You may use the standard binomial
and geometric mean and variance formulas without proof.
Geometric waiting times count trials \emph{including the first success}.
The suggested times leave 4 minutes to check your work.

\problem{1}{From outcomes to moments}{14}{13}
Let $\Omega=\{a,b,c,d\}$, with every subset an event. The outcome probabilities
and the values of a random variable $X$ are given below; $q$ is a constant.
\begin{center}
\begin{tabular}{@{}lcccc@{}}
\toprule
Outcome $\omega$ & $a$ & $b$ & $c$ & $d$\\
\midrule
$\Prob(\{\omega\})$ & $q$ & $2q$ & $q$ & $4q$\\[2pt]
$X(\omega)$ & $-2$ & $0$ & $2$ & $2$\\
\bottomrule
\end{tabular}
\end{center}
\begin{enumerate}
  \item Find $q$. List the outcomes in the event $\{X>0\}$ and find
  $\Prob(X>0)$. \partpts{3}
  \begin{solution}[.55in]
    The total mass is $8q=1$, so $q=1/8$; all masses are nonnegative.
    The event is $\{c,d\}$ and its probability is $q+4q=5/8$.
    \credit{Normalization and $q$: 1; event: 1; probability: 1.}
  \end{solution}

  \item Give the probability mass function of $X$: list each distinct
  possible value and its probability. \partpts{3}
  \begin{solution}[.65in]
    $\Prob(X=-2)=1/8$, $\Prob(X=0)=1/4$, and
    $\Prob(X=2)=5/8$, with probability zero at all other values.
    The outcomes $c,d$ contribute to the same value, so their masses add.
    \credit{Each value with its correct probability: 1.}
  \end{solution}

  \item Compute $\E[X]$, $\E[X^2]$, and $\Var(X)$. \partpts{5}
  \begin{solution}[1.05in]
    \[
      \E[X]=(-2)\frac18+0\frac14+2\frac58=1,
      \qquad \E[X^2]=4\frac18+0\frac14+4\frac58=3.
    \]
    Hence $\Var(X)=\E[X^2]-(\E[X])^2=3-1^2=2$.
    \credit{Weighted first moment and value: 2; second moment: 1;
    variance identity: 1; variance value: 1.}
  \end{solution}

  \item Let $Y=3-2X$. Find $\E[Y]$ and $\Var(Y)$, stating the
  transformation rules you use. \partpts{3}
  \begin{solution}[.7in]
    $\E[Y]=3-2\E[X]=1$ and
    $\Var(Y)=(-2)^2\Var(X)=4\cdot2=8$.
    Adding a constant leaves variance unchanged; scaling squares the factor.
    \credit{Mean with linearity: 1; squared scaling rule: 1;
    variance value: 1. A direct distribution calculation also earns full credit.}
  \end{solution}
\end{enumerate}

\newpage
\problem{2}{Bernoulli indicators}{10}{9}
Roll a fair six-sided die once. Let $I=1$ if the result is $5$ or $6$,
and let $I=0$ otherwise.
\begin{enumerate}
  \item State the distribution of $I$, including its parameter, and give
  $\Prob(I=0)$ and $\Prob(I=1)$. \partpts{2}
  \begin{solution}[1.25in]
    $I\sim\operatorname{Bernoulli}(1/3)$, with
    $\Prob(I=0)=4/6=2/3$ and $\Prob(I=1)=2/6=1/3$.
    \credit{Bernoulli family and parameter: 1; both masses: 1.}
  \end{solution}

  \item Use the two possible values of $I$ to calculate $\E[I]$,
  $\E[I^2]$, and $\Var(I)$. Show the weighted sums rather than only
  quoting a Bernoulli formula. \partpts{4}
  \begin{solution}[2in]
    \[
      \E[I]=0\cdot\frac23+1\cdot\frac13=\frac13,
      \qquad
      \E[I^2]=0^2\cdot\frac23+1^2\cdot\frac13=\frac13.
    \]
    Therefore $\Var(I)=1/3-(1/3)^2=2/9$.
    The identity $I^2=I$ also explains why the first two moments agree.
    \credit{First-moment sum: 1; second-moment sum or $I^2=I$: 1;
    variance identity: 1; value: 1.}
  \end{solution}

  \item On the \emph{same roll}, let $J=1-I$. Find $\E[I+J]$ and
  $\Var(I+J)$. A student claims
  $\Var(I+J)=\Var(I)+\Var(J)=4/9$.
  Explain why this reasoning fails here. \partpts{4}
  \begin{solution}[2.05in]
    Since $I+J=1$ on every outcome, $\E[I+J]=1$ and
    $\Var(I+J)=0$. Linearity of expectation is valid without independence.
    However, $I$ and $J$ are dependent: knowing either determines the other.
    Thus the independent-variance addition rule does not apply.
    For example, $\Prob(I=1,J=1)=0$, whereas
    $\Prob(I=1)\Prob(J=1)=(1/3)(2/3)=2/9$.
    \credit{Mean: 1; variance: 1; identify dependence and explain why
    it invalidates this use of variance addition: 2. A covariance
    calculation is not required.}
  \end{solution}
\end{enumerate}

\newpage
\problem{3}{Binomial counts}{12}{11}
A bag contains one red ball and three blue balls. Draw one ball uniformly
at random, record its color, and replace it. Repeat for a total of six
draws. Assume the draws are independent. Let $N$ be the number of red
balls recorded.
\begin{enumerate}
  \item State the distribution of $N$, its parameters, and its possible
  values. Explain why the model applies. \partpts{3}
  \begin{solution}[1.2in]
    $N\sim\operatorname{Bin}(6,1/4)$, supported on $0,1,\ldots,6$.
    There are six fixed trials, each records red or not red, the red
    probability stays $1/4$ because of replacement, and the draws are
    independent by assumption.
    \credit{Family and parameters: 1; support: 1; modeling assumptions: 1.}
  \end{solution}

  \item Find $\Prob(N=2)$. Explain the counting factor in your expression.
  \partpts{3}
  \begin{solution}[1.5in]
    \[
      \Prob(N=2)=\binom62\left(\frac14\right)^2
                            \left(\frac34\right)^4
                    =\frac{1215}{4096}.
    \]
    The factor $\binom62=15$ chooses the two red positions among the six
    draws. Each such red/blue sequence has the same probability.
    \credit{Correct probability of a specified sequence: 1;
    binomial factor: 1; explanation of that factor: 1.}
  \end{solution}

  \item Find the probability that at least one red ball is recorded.
  \partpts{2}
  \begin{solution}[1.25in]
    \[
      \Prob(N\ge1)=1-\Prob(N=0)
        =1-\left(\frac34\right)^6=\frac{3367}{4096}.
    \]
    \credit{Complement or equivalent correct sum: 1;
    correct probability expression: 1.}
  \end{solution}

  \item Compute $\E[N]$ and $\Var(N)$. If $N$ is written as a sum of
  six success indicators, why may their variances be added? \partpts{4}
  \begin{solution}[1.65in]
    \[
      \E[N]=6\cdot\frac14=\frac32,
      \qquad \Var(N)=6\cdot\frac14\cdot\frac34=\frac98.
    \]
    Each indicator records one of the independent draws, so the indicators
    are independent and their variances add. Their expectations add by
    linearity, which itself does not require independence.
    \credit{Mean: 1; variance formula: 1; variance value: 1;
    independence justification: 1.}
  \end{solution}
\end{enumerate}

\newpage
\problem{4}{Geometric waiting times}{14}{13}
A device attempts to send a message repeatedly. Each attempt succeeds
with probability $1/4$, independently of all other attempts. Let $T$ be
the number of attempts \emph{up to and including the first success}.
\begin{enumerate}
  \item Give the possible values of $T$ and a formula for
  $\Prob(T=k)$. Explain the formula and use a geometric series to check
  that its probabilities sum to $1$. \partpts{4}
  \begin{solution}[1.1in]
    The support is $1,2,3,\ldots$, and
    $\Prob(T=k)=(3/4)^{k-1}(1/4)$: the first $k-1$ attempts fail and
    the $k$th succeeds. Independence lets us multiply these probabilities.
    \[
      \sum_{k=1}^{\infty}\Prob(T=k)
        =\frac14\sum_{j=0}^{\infty}\left(\frac34\right)^j
        =\frac{1/4}{1-3/4}=1.
    \]
    \credit{Support: 1; PMF: 1; failure/success explanation: 1;
    normalization: 1.}
  \end{solution}

  \item Find $\Prob(T=3)$ and $\Prob(T>3)$. \partpts{2}
  \begin{solution}[.75in]
    $\Prob(T=3)=(3/4)^2(1/4)=9/64$.
    The event $T>3$ means that the first three attempts fail, so
    $\Prob(T>3)=(3/4)^3=27/64$.
    \credit{Each correct probability: 1.}
  \end{solution}

  \item Find $\E[T]$ and $\Var(T)$. \partpts{3}
  \begin{solution}[.8in]
    With $p=1/4$ and trials counted through the first success,
    \[
      \E[T]=\frac1p=4,\qquad
      \Var(T)=\frac{1-p}{p^2}=\frac{3/4}{(1/4)^2}=12.
    \]
    \credit{Mean: 1; variance formula: 1; variance value: 1.}
  \end{solution}

  \item Let $F$ be the number of \emph{failed} attempts before the first
  success. Express $F$ in terms of $T$, then find $\E[F]$ and $\Var(F)$.
  \partpts{2}
  \begin{solution}[.65in]
    $F=T-1$, so $\E[F]=\E[T]-1=3$ and
    $\Var(F)=\Var(T)=12$; subtracting a constant does not change variance.
    \credit{Relation and mean: 1; unchanged variance: 1.}
  \end{solution}

  \item Given that the first two attempts have failed, find the probability
  that more than five attempts in total will be needed:
  $\Prob(T>5\mid T>2)$. Justify your answer. \partpts{3}
  \begin{solution}[1.05in]
    Since $\{T>5\}\subseteq\{T>2\}$,
    \[
      \Prob(T>5\mid T>2)=\frac{\Prob(T>5)}{\Prob(T>2)}
        =\frac{(3/4)^5}{(3/4)^2}
        =\left(\frac34\right)^3=\frac{27}{64}.
    \]
    Equivalently, attempts 3, 4, and 5 must all fail; independence keeps
    their conditional failure probability at $3/4$ each. This is the
    geometric memoryless property.
    \credit{Conditional-probability setup or identify three further
    failures: 1; justified cancellation or independence: 1; result: 1.}
  \end{solution}
\end{enumerate}

\label{quiz-last-page}
\end{document}
